\documentclass[10pt]{extarticle}
\usepackage[a4paper,margin=1.5cm]{geometry}
\usepackage[T1]{fontenc}
\usepackage{amsmath,amssymb}
\usepackage{graphicx}
\usepackage[export]{adjustbox}
\usepackage{multicol}
\usepackage{enumitem}
\usepackage{booktabs,array,tabularx}
\setlength{\columnsep}{0.7cm}
\setlength{\parindent}{0pt}
\setlength{\parskip}{3pt}
\pagestyle{empty}
\makeatletter
% figure registry (figdefs.tex is generated by the extract step)
\newcommand{\qfigdef}[4]{\@namedef{qf@file@#1}{#2}\@namedef{qf@w@#1}{#3}\@namedef{qf@wide@#1}{#4}}
\newcommand{\qf@check}[1]{\@ifundefined{qf@file@#1}{\PackageError{score2}{Unknown figure id #1}{}}{}}
\newcommand{\qf@img}[4][]{\includegraphics[width=#2\dimexpr\@nameuse{qf@w@#4}pt\relax,max width=#3,max height=0.45\textheight,#1]{\@nameuse{qf@file@#4}}}
% Full-width questions: finalize wraps every question that contains a wide figure in \qfullbegin...\qfullend,
% so the WHOLE question (stem, figures, options) is set across the page in reading order, then the two columns
% resume. The question's head - its number, stem and first figure (\qfullsplit marks the end of the head) -
% is kept on one page, so a number is never left at a page foot with its figure overleaf; the rest may break
% across pages normally, so no big gaps are left. Without a \qfullsplit, the whole question is kept together
% if it is under 45% of a page. Never write these yourself.
\newif\ifqf@full
\newif\ifqf@boxed
\newsavebox{\qf@box}
\newcommand{\qf@open}{\global\qf@fulltrue\par\global\qf@boxedtrue
  \setbox\qf@box\vbox\bgroup\hsize\textwidth\linewidth\textwidth}
\newcommand{\qf@emit}[1]{\par\egroup\global\qf@boxedfalse
  \ifdim\dimexpr\ht\qf@box+\dp\qf@box\relax<#1\textheight\noindent\usebox{\qf@box}\par\else\unvbox\qf@box\fi}
\newcommand{\qfullsplit}{\ifqf@boxed\qf@emit{0.85}\fi}
\newcommand{\qf@close}{\ifqf@boxed\qf@emit{0.45}\fi\global\qf@fullfalse}
\newcommand{\qfullbegin}{\end{multicols}\qf@open}
\newcommand{\qfullend}{\qf@close\begin{multicols}{2}}
\newcommand{\qfullstart}{\qf@open}
\newcommand{\qfulldone}{\qf@close}
% \qfig[scale]{id}: block figure, centred, directly where it is written (size set by the tool).
% Fallback if a wide figure is not inside a full-width question: break out of the columns in place - no floats.
\newcommand{\qfig}[2][1]{\qf@check{#2}%
  \ifqf@full\qf@blk{#1}{#2}\else
  \ifnum\@nameuse{qf@wide@#2}=1 \expandafter\@firstoftwo\else\expandafter\@secondoftwo\fi
  {\end{multicols}\begin{center}\qf@img{#1}{\textwidth}{#2}\end{center}\begin{multicols}{2}}%
  {\qf@blk{#1}{#2}}\fi}
\newcommand{\qf@blk}[2]{\par\smallskip{\centering\qf@img{#1}{\linewidth}{#2}\par}\smallskip}
% \qfigi[scale]{id}: inline figure (inside option lists / tables / minipages); never breaks columns.
\newcommand{\qfigi}[2][1]{\qf@check{#2}\qf@img[valign=c]{#1}{\linewidth}{#2}}
\makeatother
% structure macros - lint reads these to check completeness and figure placement
\newcommand{\papersec}[1]{\par\bigskip{\centering\Large\bfseries PAPER-#1\par}\medskip}
\newcommand{\partsec}[2]{\par\medskip\noindent{\large\bfseries PART-#1: #2}\par\smallskip}
\newcommand{\sect}[1]{\par\smallskip\noindent\textbf{#1}\par}
\newcommand{\qn}[1]{\par\smallskip\noindent\textbf{#1.}~}
\newcommand{\paraq}[2]{\par\smallskip\noindent\textbf{Paragraph for Questions #1 and #2}\par}
\newenvironment{qoptions}{\begin{itemize}[leftmargin=1.8em,itemsep=1pt,topsep=2pt,label={}]}{\end{itemize}}
\newcommand{\opt}[2]{\item[(#1)] #2}
% figure registry (generated by score2_tools.py finalize)
\qfigdef{p1_phy_s1_q01_1}{Score Advanced - Test 39 (20 Apr 2025)-img/p1_phy_s1_q01_1.png}{79.0}{0}
\qfigdef{p1_phy_s1_q02_1}{Score Advanced - Test 39 (20 Apr 2025)-img/p1_phy_s1_q02_1.png}{310.1}{1}
\qfigdef{p1_phy_s1_q04_1}{Score Advanced - Test 39 (20 Apr 2025)-img/p1_phy_s1_q04_1.png}{165.4}{0}
\qfigdef{p1_phy_s2_q05_1}{Score Advanced - Test 39 (20 Apr 2025)-img/p1_phy_s2_q05_1.png}{109.1}{0}
\qfigdef{p1_phy_s2_q05_2}{Score Advanced - Test 39 (20 Apr 2025)-img/p1_phy_s2_q05_2.png}{70.7}{0}
\qfigdef{p1_phy_s2_q06_1}{Score Advanced - Test 39 (20 Apr 2025)-img/p1_phy_s2_q06_1.png}{129.2}{0}
\qfigdef{p1_phy_s2_q07_1}{Score Advanced - Test 39 (20 Apr 2025)-img/p1_phy_s2_q07_1.png}{107.1}{0}
\qfigdef{p1_phy_s2_q08_1}{Score Advanced - Test 39 (20 Apr 2025)-img/p1_phy_s2_q08_1.png}{120.5}{0}
\qfigdef{p1_phy_s2_q09_1}{Score Advanced - Test 39 (20 Apr 2025)-img/p1_phy_s2_q09_1.png}{101.5}{0}
\qfigdef{p1_phy_s2_q10_1}{Score Advanced - Test 39 (20 Apr 2025)-img/p1_phy_s2_q10_1.png}{72.7}{0}
\qfigdef{p1_phy_s3_para1_2_1}{Score Advanced - Test 39 (20 Apr 2025)-img/p1_phy_s3_para1_2_1.png}{128.1}{0}
\qfigdef{p1_phy_s3_q06_1}{Score Advanced - Test 39 (20 Apr 2025)-img/p1_phy_s3_q06_1.png}{90.5}{0}
\qfigdef{p1_chem_s1_q01_A}{Score Advanced - Test 39 (20 Apr 2025)-img/p1_chem_s1_q01_A.png}{142.4}{0}
\qfigdef{p1_chem_s1_q01_B}{Score Advanced - Test 39 (20 Apr 2025)-img/p1_chem_s1_q01_B.png}{137.6}{0}
\qfigdef{p1_chem_s1_q01_C}{Score Advanced - Test 39 (20 Apr 2025)-img/p1_chem_s1_q01_C.png}{101.5}{0}
\qfigdef{p1_chem_s1_q01_D}{Score Advanced - Test 39 (20 Apr 2025)-img/p1_chem_s1_q01_D.png}{141.7}{0}
\qfigdef{p1_chem_s1_q03_1}{Score Advanced - Test 39 (20 Apr 2025)-img/p1_chem_s1_q03_1.png}{64.2}{0}
\qfigdef{p1_chem_s1_q03_2}{Score Advanced - Test 39 (20 Apr 2025)-img/p1_chem_s1_q03_2.png}{49.4}{0}
\qfigdef{p1_chem_s1_q03_A}{Score Advanced - Test 39 (20 Apr 2025)-img/p1_chem_s1_q03_A.png}{43.8}{0}
\qfigdef{p1_chem_s1_q03_B}{Score Advanced - Test 39 (20 Apr 2025)-img/p1_chem_s1_q03_B.png}{44.9}{0}
\qfigdef{p1_chem_s1_q03_C}{Score Advanced - Test 39 (20 Apr 2025)-img/p1_chem_s1_q03_C.png}{43.2}{0}
\qfigdef{p1_chem_s1_q03_D}{Score Advanced - Test 39 (20 Apr 2025)-img/p1_chem_s1_q03_D.png}{43.2}{0}
\qfigdef{p1_chem_s2_q07_C}{Score Advanced - Test 39 (20 Apr 2025)-img/p1_chem_s2_q07_C.png}{49.9}{0}
\qfigdef{p1_chem_s2_q07_D}{Score Advanced - Test 39 (20 Apr 2025)-img/p1_chem_s2_q07_D.png}{46.6}{0}
\qfigdef{p1_chem_s3_q02_1}{Score Advanced - Test 39 (20 Apr 2025)-img/p1_chem_s3_q02_1.png}{86.3}{0}
\qfigdef{p1_chem_s3_para3_4_1}{Score Advanced - Test 39 (20 Apr 2025)-img/p1_chem_s3_para3_4_1.png}{243.9}{0}
\qfigdef{p1_chem_s3_para5_6_1}{Score Advanced - Test 39 (20 Apr 2025)-img/p1_chem_s3_para5_6_1.png}{103.1}{0}
\qfigdef{p2_phy_s1_q02_1}{Score Advanced - Test 39 (20 Apr 2025)-img/p2_phy_s1_q02_1.png}{202.1}{0}
\qfigdef{p2_phy_s1_q03_1}{Score Advanced - Test 39 (20 Apr 2025)-img/p2_phy_s1_q03_1.png}{159.5}{0}
\qfigdef{p2_phy_s1_q04_1}{Score Advanced - Test 39 (20 Apr 2025)-img/p2_phy_s1_q04_1.png}{107.3}{0}
\qfigdef{p2_phy_s1_q06_1}{Score Advanced - Test 39 (20 Apr 2025)-img/p2_phy_s1_q06_1.png}{98.8}{0}
\qfigdef{p2_phy_s2_para7_8_1}{Score Advanced - Test 39 (20 Apr 2025)-img/p2_phy_s2_para7_8_1.png}{211.4}{0}
\qfigdef{p2_phy_s2_para7_8_2}{Score Advanced - Test 39 (20 Apr 2025)-img/p2_phy_s2_para7_8_2.png}{210.9}{0}
\qfigdef{p2_phy_s2_q10_1}{Score Advanced - Test 39 (20 Apr 2025)-img/p2_phy_s2_q10_1.png}{79.9}{0}
\qfigdef{p2_phy_s3_para1_2_1}{Score Advanced - Test 39 (20 Apr 2025)-img/p2_phy_s3_para1_2_1.png}{76.0}{0}
\qfigdef{p2_phy_s3_q04_1}{Score Advanced - Test 39 (20 Apr 2025)-img/p2_phy_s3_q04_1.png}{74.3}{0}
\qfigdef{p2_phy_s3_para5_6_1}{Score Advanced - Test 39 (20 Apr 2025)-img/p2_phy_s3_para5_6_1.png}{65.4}{0}
\qfigdef{p2_phy_s4_q07_1}{Score Advanced - Test 39 (20 Apr 2025)-img/p2_phy_s4_q07_1.png}{78.4}{0}
\qfigdef{p2_phy_s4_q08_1}{Score Advanced - Test 39 (20 Apr 2025)-img/p2_phy_s4_q08_1.png}{169.4}{0}
\qfigdef{p2_chem_s1_q01_A}{Score Advanced - Test 39 (20 Apr 2025)-img/p2_chem_s1_q01_A.png}{136.1}{0}
\qfigdef{p2_chem_s1_q01_2}{Score Advanced - Test 39 (20 Apr 2025)-img/p2_chem_s1_q01_2.png}{24.2}{0}
\qfigdef{p2_chem_s1_q01_B}{Score Advanced - Test 39 (20 Apr 2025)-img/p2_chem_s1_q01_B.png}{134.5}{0}
\qfigdef{p2_chem_s1_q01_4}{Score Advanced - Test 39 (20 Apr 2025)-img/p2_chem_s1_q01_4.png}{27.9}{0}
\qfigdef{p2_chem_s1_q01_5}{Score Advanced - Test 39 (20 Apr 2025)-img/p2_chem_s1_q01_5.png}{25.4}{0}
\qfigdef{p2_chem_s1_q01_6}{Score Advanced - Test 39 (20 Apr 2025)-img/p2_chem_s1_q01_6.png}{37.4}{0}
\qfigdef{p2_chem_s1_q01_D}{Score Advanced - Test 39 (20 Apr 2025)-img/p2_chem_s1_q01_D.png}{216.2}{0}
\qfigdef{p2_chem_s1_q02_1}{Score Advanced - Test 39 (20 Apr 2025)-img/p2_chem_s1_q02_1.png}{281.0}{1}
\qfigdef{p2_chem_s1_q02_A}{Score Advanced - Test 39 (20 Apr 2025)-img/p2_chem_s1_q02_A.png}{60.5}{0}
\qfigdef{p2_chem_s1_q02_B}{Score Advanced - Test 39 (20 Apr 2025)-img/p2_chem_s1_q02_B.png}{64.4}{0}
\qfigdef{p2_chem_s1_q02_C}{Score Advanced - Test 39 (20 Apr 2025)-img/p2_chem_s1_q02_C.png}{58.6}{0}
\qfigdef{p2_chem_s1_q02_D}{Score Advanced - Test 39 (20 Apr 2025)-img/p2_chem_s1_q02_D.png}{66.2}{0}
\qfigdef{p2_chem_s1_q04_1}{Score Advanced - Test 39 (20 Apr 2025)-img/p2_chem_s1_q04_1.png}{203.8}{0}
\qfigdef{p2_chem_s2_para9_10_1}{Score Advanced - Test 39 (20 Apr 2025)-img/p2_chem_s2_para9_10_1.png}{122.1}{0}
\qfigdef{p2_chem_s3_para5_6_1}{Score Advanced - Test 39 (20 Apr 2025)-img/p2_chem_s3_para5_6_1.png}{193.6}{0}
\qfigdef{p2_chem_s4_q07_1}{Score Advanced - Test 39 (20 Apr 2025)-img/p2_chem_s4_q07_1.png}{50.7}{0}
\raggedcolumns
\makeatletter\@beginparpenalty=10000\makeatother
\newcolumntype{L}[1]{>{\raggedright\arraybackslash}p{#1}}
\begin{document}
\begin{center}{\Large\bfseries Score Advanced --- Test 39 (20 Apr 2025)}\end{center}
\begin{multicols}{2}
\raggedcolumns
\papersec{1}
\partsec{1}{PHYSICS}
\sect{SECTION-I(i)}
\qn{1} An electric circuit consists of three charged capacitors. The capacitances of the capacitors and the values of their charges are given in the diagram. At some moment, the plates of capacitor X are released, allowing them to move freely toward each other. As a result of inelastic collisions, the plates eventually touch. What amount of heat is released in the system as a result?\par
\qfig{p1_phy_s1_q01_1}
\begin{qoptions}
\opt{A}{$\dfrac{Q^2}{2C}$}
\opt{B}{$\dfrac23\dfrac{Q^2}{C}$}
\opt{C}{$\dfrac{Q^2}{C}$}
\opt{D}{$\dfrac25\dfrac{Q^2}{C}$}
\end{qoptions}
\qfullbegin
\qn{2} Match the curve in figure-1 to the curve in figure-2 having the same quality factor in LCR series circuit with AC voltage. Here $\phi$ is phase difference between current and voltage.\par
\qfig{p1_phy_s1_q02_1}
\qfullsplit
\begin{qoptions}
\opt{A}{1 $\to$ a, 2 $\to$ b, 3 $\to$ c, 4 $\to$ d}
\opt{B}{1 $\to$ d, 2 $\to$ c, 3 $\to$ b, 4 $\to$ a}
\opt{C}{1 $\to$ a, 2 $\to$ b, 3 $\to$ d, 4 $\to$ c}
\opt{D}{1 $\to$ d, 2 $\to$ c, 3 $\to$ c, 4 $\to$ d}
\end{qoptions}
\qfullend
\qn{3} A circular table has radius R and N $>$ 2 equally spaced legs of length h attached to its perimeter. Suppose the table has a uniform mass density with total mass m. Assuming the table does not slip, the minimum horizontal force needed to tip over the table is :-
\begin{qoptions}
\opt{A}{$\dfrac{\mathrm{mgR}}{\mathrm h}$}
\opt{B}{$\dfrac{\mathrm{mgR}}{\mathrm h}\cos\!\left(\dfrac{\pi}{\mathrm N}\right)$}
\opt{C}{$\dfrac{\mathrm{mgR}}{\mathrm h}\tan\!\left(\dfrac{\mathrm N-2}{2\mathrm N}\pi\right)$}
\opt{D}{$\dfrac{\mathrm{mgR}}{\mathrm h}\sin\!\left(\dfrac{\pi}{2\mathrm N}\right)$}
\end{qoptions}
\qn{4} The schematic provided has all resistors with a resistance of three ohms, and the battery voltage is 48 volts. All ammeters are ideal.\par
Statement-1 : Readings of all ammeters is same.\par
Statement-2 : Reading of ammeter 1 is 5A.\par
Statement-3 : Reading of ammeter 2 is 7A.\par
Statement-4 : Reading of ammeter 3 is 7A.\par
\qfig{p1_phy_s1_q04_1}
\begin{qoptions}
\opt{A}{All statements are correct.}
\opt{B}{Only statement-1 is correct.}
\opt{C}{Statement-1 \& statement-4 are correct.}
\opt{D}{Statement-2 \& statement-3 are correct.}
\end{qoptions}
\sect{SECTION-I(ii)}
\qn{5} A particle is constrained to move on the inner surface of a frictionless parabolic bowl whose cross-section has equation z = kr$^2$. The particle begins at a height z$_0$ above the bottom of the bowl with a horizontal velocity v$_0$ along the surface of the bowl. The acceleration due to gravity is g.\par
\qfig{p1_phy_s2_q05_1}
\qfig{p1_phy_s2_q05_2}
\begin{qoptions}
\opt{A}{The value of v$_0$ for which the particle moves in a horizontal circle is $\sqrt{\mathrm{gz}_0}$.}
\opt{B}{For v$_0 > \sqrt{2\mathrm{gz}_0}$, the maximum height reached by the particle is $\dfrac{\mathrm v_0^2}{2\mathrm g}$}
\opt{C}{If v$_0$ = 0 and z$_0$ is very small then the time in which the particle returns to its initial point is $4\sqrt{\dfrac{2\mathrm z_0}{\mathrm g}}$.}
\opt{D}{If v$_0$ = 0 and z$_0$ is very small, then the time in which the particle returns to its initial point is $\dfrac{\pi\sqrt2}{\sqrt{\mathrm{kg}}}$.}
\end{qoptions}
\qn{6} A cart for studying uniformly accelerated motion is driven by a weight, as shown in Fig. A descending weight of mass m rotates a rear axle via a rope passing over a pulley, with the rope sufficiently wound around the axle. Initially, the weight is at a height h above the cart. When released, the cart starts moving. Each of the four wheels of the cart has a mass m$_1$ and a radius r$_1$. The rear and front axles have a mass m$_2$ and a radius r$_2$. The wheels and axles are made as solid homogeneous cylinders. The motion of the weight is guided by a smooth vertical tube. The mass of the cart without the weight, axles, and wheels is negligible. The mass of the pulley, friction of the weight inside the tube, and rolling resistance of the wheels are considered negligible.\par
\qfig{p1_phy_s2_q06_1}
\begin{qoptions}
\opt{A}{The acceleration of the cart is\par {\small $\dfrac{\mathrm{mgr}_2}{\mathrm r_1\left[\mathrm m\left(1+\frac{\mathrm r_2^2}{\mathrm r_1^2}\right)+6\mathrm m_1+\mathrm m_2\left(2+\frac{\mathrm r_2^2}{\mathrm r_1^2}\right)\right]}$}}
\opt{B}{The speed of the cart till the weight falls on the cart is $\sqrt{\dfrac{2\mathrm{mgh}}{\mathrm m\left(1+\frac{\mathrm r_2^2}{\mathrm r_1^2}\right)+6\mathrm m_1+\mathrm m_2\left(2+\frac{\mathrm r_2^2}{\mathrm r_1^2}\right)}}$}
\opt{C}{The tension in the string while the weight is falling is $\mathrm m\left(\mathrm g-\mathrm a\dfrac{\mathrm r_2}{\mathrm r_1}\right)$, where a is the acceleration of the cart.}
\opt{D}{The tension in the string while the weight is falling is $\mathrm m\left(\mathrm g-\mathrm a\dfrac{\mathrm r_2^2}{\mathrm r_1^2}\right)$, where a is the acceleration of the cart.}
\end{qoptions}
\qn{7} In the circuit shown, the switch S has been closed for a long time. At time t = 0, the switch is opened. It remains open for a time T, after which it is closed again.\par
\qfig{p1_phy_s2_q07_1}
\begin{qoptions}
\opt{A}{Voltage drop across 100 k$\Omega$ resister is 10e$^{-t/1.5}$ V for t $<$ T.}
\opt{B}{Voltage drop across 100 k$\Omega$ resister is $\dfrac{20}3$e$^{-t/1.5}$ V for t $<$ T.}
\opt{C}{Voltage drop across 100 k$\Omega$ resister is 10(e$^{\mathrm T-t}-$e$^{(\mathrm T-3t)/3}$)V for t $>$ T.}
\opt{D}{The time constant of the circuit is 1.5 sec for t $>$ T.}
\end{qoptions}
\qn{8} One mole of ideal monoatomic gas follows cyclic process as shown in the figure. Choose the correct option(s) :\par
\qfig{p1_phy_s2_q08_1}
\begin{qoptions}
\opt{A}{$\dfrac{\mathrm T_B+\mathrm T_C}2$ = 4T$_0$}
\opt{B}{Molar heat capacity for process B $\to$ A is 2R}
\opt{C}{Molar heat capacity for process D $\to$ C is 4R}
\opt{D}{Work done by gas during D to C is $\dfrac{3\mathrm{RT}_0}2$}
\end{qoptions}
\qn{9} A cubical container has 2 fixed vertical walls and two vertical walls can freely rotate about one of the horizontal edges. The right wall can rotate about upper edge and the left wall can rotate about the bottom edge. Water is filled upto the top and forces F$_1$ and F$_2$ (perpendicular to the walls) are applied to keep the walls at from opening as shown in figure. Given that the cube is at rest on a rough horizontal surface. (Ignore atmospheric pressure)\par
\qfig{p1_phy_s2_q09_1}
\begin{qoptions}
\opt{A}{Friction on the box from the surfaces is to the left.}
\opt{B}{Friction on the box from the surfaces is to the right.}
\opt{C}{Normal reaction is shifted towards left wall.}
\opt{D}{Normal reaction passes through the center.}
\end{qoptions}
\qn{10} At the center of a tube with radius R, a source emits charged particles (see diagram). Each particle has a charge q $>$ 0 and mass m. All particles are emitted from the source in the plane of the drawing with the same speed v, evenly distributed in the upper half of the tube. A detector of length $\ell$, which registers the total number of particles that reach it in one unit of time, is fixed on the surface of the tube as shown in the diagram. Along the entire tube (in the direction ``toward us''), a uniform magnetic field of induction B is applied. For what values of B will the detector readings be at their maximum possible values? The interaction between charged particles and radiation is negligible. [Take $\ell=\dfrac{\pi\mathrm R}6$]\par
\qfig{p1_phy_s2_q10_1}
\begin{qoptions}
\opt{A}{$\dfrac{\mathrm{mv}}{\mathrm{qR}}$}
\opt{B}{$\dfrac32\dfrac{\mathrm{mv}}{\mathrm{qR}}$}
\opt{C}{$\dfrac{2\mathrm{mv}}{\mathrm{qR}}$}
\opt{D}{$\dfrac52\dfrac{\mathrm{mv}}{\mathrm{qR}}$}
\end{qoptions}
\sect{SECTION-II(i)}
\paraq{1}{2}
Figure shows the circuit of a potentiometer. The length of the potentiometer wire AB is 50 cm. The emf of the battery is 4 volt, having negligible internal resistance. Values of resistances R$_1$ and R$_2$ are 15 ohm and 5 ohm respectively. When both the keys are open, the null point is obtained at a distance of 31.25 cm from end A. Given resistance of wire AB = 10 $\Omega$.\par
\qfig{p1_phy_s3_para1_2_1}
\qn{1} Potential gradient of wire AB at null point when K$_2$ is open and K$_1$ is open is \underline{\hspace{1.5cm}} V/cm
\qn{2} The balance length when key K$_2$ is open and K$_1$ is closed is given by \underline{\hspace{1.5cm}} cm.
\paraq{3}{4}
The blue of the sky and the red of the sunset are due to a phenomenon called ``scattering''. When sunlight passes through the earth's atmosphere, much of the light is picked up by the air molecules and given out again in some other direction. Experiments show, in agreement with the theory of scattering, that the shortest waves are scattered more readily then longer waves. To more specific, the scattering is inversely proportional to the fourth power of the wavelength.\par
scattering $\propto\dfrac1{\lambda^4}$\par
For every ten violet wave ($\lambda$ = 4 $\times$ 10$^{-7}$ m) scattered from a beam, there is only one red wave ($\lambda$ = 7 $\times$ 10$^{-7}$ m). Scattering of different colours will be in the ratio as given below :\par
{\footnotesize\setlength{\tabcolsep}{4pt}\noindent\begin{tabular}{|c|c|c|c|c|c|}\hline
Violet & Blue & Green & Yellow & Orange & Red \\ \hline
10 & 6 & 3 & 2.5 & 2 & 1 \\ \hline
\end{tabular}}\par
\qn{3} Wavelength of yellow light will be approximately $\alpha$ $\times$ 10$^{-7}$ m. Write $\alpha$.
\qn{4} Scattering of wavelength corresponding to threshold wavelength of photoelectric setup of work function 6.2 eV will be nearly $\eta$ times that violet colour. The value of $\eta$ is (Give hC = 12400 eV \AA)
\paraq{5}{6}
Two particles with positive charge e and mass M, and a third particle with negative charge $-$ e and mass m = $\dfrac M2$, move in circular orbits, maintaining the configuration of an isosceles triangle (see the figure). The planes of the particles' orbits are perpendicular to the base of the triangle.
\qn{5} Find the ratio $\dfrac{\mathrm R_-}{\mathrm R_+}$ where R$_-$ is the radius of the orbit of the negatively charged particle, and R$_+$ is the radius of the orbit of the positively charged particles. Round your answer to the nearest whole number.
\qn{6} Find the ratio of the base length L of the triangle to the length of its side R. Round your answer to two decimal places.\par
\qfig{p1_phy_s3_q06_1}
\sect{SECTION-II(ii)}
\qn{7} A steel girder is planted securely between two sides of a ravine in order to provide a bridge. The total cross-sectional area of the girder is 30 cm$^2$, and the length of the girder is 4.0 m. Installed at a temperature of 5$^\circ$C, the temperature now rises to 20 $^\circ$C. Calculate the force exerted (in N) by the girder due to the change in temperature, assuming the ends do not move. Young modulus of steel = 2.0 $\times$ 10$^{11}$ Pa. Coefficient of linear expansion is 1.2 $\times$ 10$^{-7}$ $^\circ$C$^{-1}$.
\qn{8} A closed cylindrical soft drink can of height 12cm and radius 4cm is filled up to height 9cm. Volume above liquid has air at a pressure 10 kPa. Soft drink is incompressible and has specific gravity 0.8. Taken far away from all celestial bodies in a spaceship where can is rotated about its axis with constant angular velocity $\dfrac{250}{\sqrt3}$ rad/s, it exerts pressure (10000 $\times$ N) Pascal on curved surface of can. N is \ldots\ldots
\qn{9} An ornament is lying at the bed of a pond of depth 8$\sqrt3$m. Pond is full of water whose refractive index is $\sqrt{\dfrac32}$. A diver at the surface of pond is trying to look at ornament such that her line of vision makes 30$^\circ$ with the surface of pond. Perpendicular distance of image of ornament from surface of water in metres is\underline{\hspace{1.5cm}}.
\raggedcolumns
\papersec{1}
\partsec{2}{CHEMISTRY}
\sect{SECTION-I(i)}
\qn{1} Which of the following reaction represents the correct major product :
\begin{qoptions}
\opt{A}{\qfigi{p1_chem_s1_q01_A}}
\opt{B}{\qfigi{p1_chem_s1_q01_B}}
\opt{C}{\qfigi{p1_chem_s1_q01_C}}
\opt{D}{\qfigi{p1_chem_s1_q01_D}}
\end{qoptions}
\qn{2} A solution of weak acid HA was titrated with strong base KOH. The weak acid is completely neutralized by adding 60 mL of 0.1 M KOH. Now, 20mL of 0.1 M HCl was added to titrated solution, the pH was found to be 5.3. What will be the pH of solution obtained by mixing 50 mL of 0.1 M KOH and 50 mL of 0.1 M HA? (log 2= 0.3)
\begin{qoptions}
\opt{A}{8.00}
\opt{B}{8.7}
\opt{C}{8.85}
\opt{D}{9.00}
\end{qoptions}
\qn{3} In the given reaction sequence the correct compound P is :\par
\qfigi{p1_chem_s1_q03_1}\qfigi{p1_chem_s1_q03_2}
\begin{qoptions}
\opt{A}{\qfigi{p1_chem_s1_q03_A}}
\opt{B}{\qfigi{p1_chem_s1_q03_B}}
\opt{C}{\qfigi{p1_chem_s1_q03_C}}
\opt{D}{\qfigi{p1_chem_s1_q03_D}}
\end{qoptions}
\qn{4} Which of the following statement is correct?
\begin{qoptions}
\opt{A}{[CoF$_6$]$^{3-}$ is d$^2$sp$^3$ hybridized and diamagnetic in nature}
\opt{B}{[NiCl$_4$]$^{2-}$ is sp$^3$ hybridized and diamagnetic in nature.}
\opt{C}{[Ni(CO)$_4$] is dsp$^2$ hybridized and diamagnetic in nature}
\opt{D}{[Cu(NH$_3$)$_4$]$^{2+}$ is dsp$^2$ hybridized and paramagnetic in nature.}
\end{qoptions}
\sect{SECTION-I(ii)}
\qn{5} Among the following, the correct statement(s) is/are :
\begin{qoptions}
\opt{A}{The two six-membered cyclic hemiacetal forms of D-(+)-glucose are called as C2 epimers.}
\opt{B}{Invert sugar is prepared by acid catalyzed hydrolysis of maltose.}
\opt{C}{DNA contains cytosine and thymine as pyrimidine bases.}
\opt{D}{Glycine is an optically inactive amino acid.}
\end{qoptions}
\qn{6} Which of the following option is/are correct for ionic solids?
\begin{qoptions}
\opt{A}{In ZnS (Zinc blende) structure, distance between two nearest Zn$^{+2}$ ions is 5\AA, then distance between two nearest S$^{2-}$ ions is 0.707\AA.}
\opt{B}{In NaCl (rock salt) structure, the second nearest cation for any anion placed at corner is present at a distance of $\dfrac{\sqrt3}2$a(where `a' = edge length of unit cell) and number of second nearest cations are 8.}
\opt{C}{In NaCl (rock salt) structure, the third nearest anion for any anion placed at corner is present at a distance of $\sqrt{\dfrac32}$ a (where `a' = edge length of unit cell).}
\opt{D}{The packing fraction of CsCl (8 : 8) structure is $\dfrac{\sqrt3\pi}4$ (approx.) (Given that : r$_{\mathrm{Cs}^+}$ = 170 pm, r$_{\mathrm{Cl}^-}$ = 180 pm)}
\end{qoptions}
\qn{7} Which of the following compound(s) will not give silver mirror on reaction with ammoniacal silver nitrate solution ?
\begin{qoptions}
\opt{A}{formic acid}
\opt{B}{Propane}
\opt{C}{\qfigi{p1_chem_s2_q07_C}}
\opt{D}{\qfigi{p1_chem_s2_q07_D}}
\end{qoptions}
\qn{8} 2.84 g of mixture of CaCO$_3$ and MgCO$_3$ was dissolved in 250 mL of 0.4 M HCl(excess) solution which was then diluted to 400 mL. If 10 mL of resulting solution requires 10 mL of $\dfrac{\mathrm M}{20}$Na$_2$CO$_3$ for neutralization. Then,
\begin{qoptions}
\opt{A}{The mass percentage of CaCO$_3$ in original mixture is 70.42 \% (approx.)}
\opt{B}{The mole percentage of MgCO$_3$ in original mixture is 33.33\% (approx.)}
\opt{C}{The number of equivalents of HCl used with original mixture was 0.06.}
\opt{D}{The number of mole of HCl neutralized by Na$_2$ CO$_3$ is 0.01.}
\end{qoptions}
\qn{9} The incorrect statement(s) related to metallurgical processes is (are)
\begin{qoptions}
\opt{A}{In leaching process of aluminium ore (Bayer process), a small amount of freshly prepared sample of hydrated Al$_2$O$_3$ is added to induce precipitation.}
\opt{B}{Wrought iron is the purest form of commercial iron and is prepared from cast iron by reducing elemental impurities in a furnace lined with hematite.}
\opt{C}{Copper matte primarily consists of Cu$_2$O and a smaller amount of FeS.}
\opt{D}{Extraction of zinc from zinc blende involves calcination followed by reduction with carbon.}
\end{qoptions}
\qn{10} The incorrect statement(s) related to colloids is/are.
\begin{qoptions}
\opt{A}{The colloidion is a 4\% solution of nitro-cellulose in a mixture of alcohol and ester.}
\opt{B}{A mixture of milk and water appears red when viewed by the reflected light and blue when viewed by the transmitted light.}
\opt{C}{Sols of starch, gum, clay and methylene blue are negatively charged colloidal sol.}
\opt{D}{When the ionic strength of a colloidal solution is increased, thickness of double layer is decreased and the colloid gets precipitated.}
\end{qoptions}
\sect{SECTION-II(i)}
\paraq{1}{2}
A compound P(C$_9$H$_{10}$) gives benzaldehyde and aliphatic aldehyde on reductive ozonolysis. This aliphatic aldehyde gives positive iodoform test. Identify P and answer the given questions.
\qn{1} What is the total number of major organic products formed [through most stable carbocation] when P is reacted with HBr/CCl$_4$ ?
\qn{2} Find degree of unsaturation of the product Q ?\par
\qfig{p1_chem_s3_q02_1}
\paraq{3}{4}
\qfig{p1_chem_s3_para3_4_1}
\qn{3} The number of unpaired electrons in compound `B' is `p' and number of atoms in a molecule of compound E are `q', then $|\mathrm p-\mathrm q|$ is:
\qn{4} When `A' is treated with conc. H$_2$SO$_4$, a gas is produced which is colorless, odourless and poisonous. The molar mass of gas (g/mol) is `x' and oxidation state of central metal in compound `A' is `y'. Then, value of $|\mathrm x+\mathrm y|$ is :
\paraq{5}{6}
0.06 moles of an ideal monoatomic gas undergoes a process as shown in figure. The temperature at point `A' is 100 K.\par
(Use : R = $\dfrac{25}3$J/mol.K, 1L.bar = 100J )\par
\qfig{p1_chem_s3_para5_6_1}
\qn{5} The magnitude of work done (in Joule) in cyclic process ABCD is
\qn{6} The magnitude of enthalpy change (in L.bar) in process CD is:
\sect{SECTION-II(ii)}
\qn{7} In the following reaction sequence, the amount of T (in g) formed from 12 mole CaC$_2$ is\underline{\hspace{0.8cm}}.\par
[Atomic weight in gmol$^{-1}$ : H = 1, C = 12, N = 14, O = 16, Ca = 40.] The yield (\%) corresponding to the product in each step is given in the parenthesis\par
{\footnotesize $\mathrm{CaC_2}\xrightarrow{\mathrm{H_2O}}\underset{(100\%)}{\mathrm P}\xrightarrow[\text{iron tube}]{\text{red hot}}\underset{(60\%)}{\mathrm Q}\xrightarrow{\text{conc. HNO}_3,\ \text{conc. H}_2\text{SO}_4}\underset{(50\%)}{\mathrm R}\xrightarrow{\mathrm{Sn/HCl}}\underset{(50\%)}{\mathrm S}\xrightarrow{\mathrm{Ac_2O}}\underset{(100\%)}{\mathrm T}+\mathrm{CH_3COOH}$}\par
\qn{8} A sample of excited single electron ions can emit maximum 10 different spectral lines when de-excited to ground level. The maximum energy emitted during transition is 117.5 eV. Now, photons of minimum energy emitted during transition strikes on a metal having work function 1.254 eV emits photoelectrons. The de-Broglie wavelength of emitted photo electron (in \AA) is:-
\qn{9} A sample of water is electrolysed during an experiment using 40 F electricity with 75\% current efficiency. The released O$_2$ is used in a bomb calorimeter to burn 3.5 g C$_2$H$_4$ gas at 300 K. The temperature of the calorimeter was found to increase from 300 K to 300.50 K due to combustion process. The magnitude of heat of combustion at constant pressure (in kJ) for 5 moles of ethene is: (Given : Heat capacity of the calorimeter is 3 kJ/K, R = 8 J/mol.K)
\raggedcolumns
\papersec{1}
\partsec{3}{MATHEMATICS}
\sect{SECTION-I(i)}
\qn{1} Straight line L: 2x + y = 4 is tangent to an ellipse whose foci are (1, 3) \& (2, 4). Straight line L$_1$ parallel to line L also touches this ellipse. Point of contact of L$_1$ to this ellipse is
\begin{qoptions}
\opt{A}{$\left(\dfrac{53}{14},\dfrac{65}{14}\right)$}
\opt{B}{$\left(\dfrac{61}{25},\dfrac{103}{25}\right)$}
\opt{C}{$\left(\dfrac{15}7,\dfrac{31}{25}\right)$}
\opt{D}{$\left(\dfrac{63}{25},\dfrac{103}{25}\right)$}
\end{qoptions}
\qn{2} $\displaystyle\int_1^2(2^x+2^{-x})\mathrm d(2^x)$ =
\begin{qoptions}
\opt{A}{$\dfrac54+\ln2$}
\opt{B}{$\dfrac32+\ln2$}
\opt{C}{$\dfrac32-\ln2$}
\opt{D}{$\dfrac78+\ln2$}
\end{qoptions}
\qn{3} $\displaystyle\lim_{x\to0}\dfrac{x\mathrm e^{\sin x}-\mathrm e^x\sin^{-1}(\sin x)}{\sin^2x-x\sin x}$ =
\begin{qoptions}
\opt{A}{$-2$}
\opt{B}{$-1$}
\opt{C}{1}
\opt{D}{2}
\end{qoptions}
\qn{4} If the standard deviation of 20 obervations x$_1$, x$_2$, x$_3$,\ldots\ldots, x$_{20}$ is 4 and that of another 20 observations y$_1$, y$_2$, y$_3$,\ldots\ldots, y$_{20}$ is 3. Also X$_\mathrm i$ = (x$_\mathrm i-\bar{\mathrm x}$)(y$_\mathrm i-\bar{\mathrm y}$) where $\bar{\mathrm x}$ and $\bar{\mathrm y}$ are means of x$_\mathrm i$'s and y$_\mathrm i$'s respectively. If $\displaystyle\sum_{\mathrm i=1}^{20}$X$_\mathrm i$ = 90, then the standard deviation of the observations x$_1-$y$_1$, x$_2-$y$_2$, x$_3-$y$_3$,\ldots\ldots, x$_{20}-$y$_{20}$ is divisible by
\begin{qoptions}
\opt{A}{5}
\opt{B}{4}
\opt{C}{3}
\opt{D}{7}
\end{qoptions}
\sect{SECTION-I(ii)}
\qn{5} $\displaystyle\int\dfrac{\mathrm x^4+32}{\mathrm x^8(\mathrm x^4-16)}\mathrm{dx}=\dfrac3\lambda\ln\left|\dfrac{\mathrm x-2}{\mathrm x+2}\right|-\dfrac\mu{256}\tan^{-1}\dfrac{\mathrm x}2+\dfrac1{\alpha\mathrm x^3}+\dfrac2{7\mathrm x^7}+\mathrm c$
\begin{qoptions}
\opt{A}{$\lambda$ = 512}
\opt{B}{$\mu$ = 5}
\opt{C}{$\mu$ = 3}
\opt{D}{$\alpha$ = 15}
\end{qoptions}
\qn{6} Let A(3, 0, 4) and B(4, 6, 0) be two given points and P($\alpha$, 0, 0) be point on x-axis such that PA + PB is minimum, then
\begin{qoptions}
\opt{A}{If point Q is on the y-axis such that QA + QB is minimum then area of $\Delta$POQ is $\dfrac{21}4$ (unit)$^2$ where O is origin}
\opt{B}{The volume of tetrahedron OPAB (where O is origin) is $\dfrac{68}5$ (unit)$^3$}
\opt{C}{Perpendicular distance from the point A to the plane containing the points P, Q, B is 4 (units)}
\opt{D}{If point Q is on the y-axis such that QA + QB is minimum then area of $\Delta$POQ is $\dfrac{17}3$ (unit)$^2$ where O is origin}
\end{qoptions}
\qn{7} If (f(x) $-$ 2)(x$^4$ + x$^2$ + 1) $-$ (f(x) + 2)(x$^2$ + x + 1)$^2$ = 0 is true for all x $\in$ R $-$ \{0\}, then which of the following statement(s) is/are true?
\begin{qoptions}
\opt{A}{f(x) has a local maxima at x = $-$1}
\opt{B}{$|$f(x)$|\geqslant$ 4, $\forall$x $\in$ R $-$ \{0\}}
\opt{C}{f(x) has a local maxima at x = 1}
\opt{D}{The value of $\displaystyle\int_{-\pi/2}^{\pi/2}$(sec x)f(x)dx = 0}
\end{qoptions}
\qn{8} If P$_\mathrm n$ denotes the number of ways in which n runners participating in a race can reach the finishing line, when one or more runners can reach at same time, then
\begin{qoptions}
\opt{A}{P$_4$ = 74}
\opt{B}{P$_4$ = 75}
\opt{C}{P$_5$ = 541}
\opt{D}{P$_5$ = 521}
\end{qoptions}
\qn{9} Let roots of quadratic equation 8x$^2-$14x + 3 = 0 are $\alpha$ and $\beta$ where $\alpha<\beta$. Which of the following equation has $\left(2\alpha+\sqrt{\dfrac\beta2}\,\mathrm i\right)^{50}$ and $\left(2\alpha-\sqrt{\dfrac\beta2}\,\mathrm i\right)^{50}$ as roots :
\begin{qoptions}
\opt{A}{x$^2$ + x + 1 = 0}
\opt{B}{x$^3$ + 1 = 0}
\opt{C}{x$^3-$ 1 = 0}
\opt{D}{x$^2-$ x + 1 = 0}
\end{qoptions}
\qn{10} If $\ln$x $\leqslant\lambda$x$^3$, $\forall$x $\in$ (0, $\infty$) then possible value/s of $\lambda$ are
\begin{qoptions}
\opt{A}{$\dfrac15$}
\opt{B}{$\dfrac1{10}$}
\opt{C}{1}
\opt{D}{$\dfrac1{\mathrm e^2}$}
\end{qoptions}
\sect{SECTION-II(i)}
\paraq{1}{2}
Let a, b and c be real numbers such that the system of linear equations\par
10x + 11y + 12z = a\par
13x + 14y + 15z = b\par
16x + 17y + 18z = c $-$ 3\par
is consistent. Let $|$x$|$ denotes the determinant of matrix\par
X = $\begin{bmatrix}\mathrm a&2&1\\\mathrm b&1&0\\\mathrm c&0&-1\end{bmatrix}$\par
Let P be the plane containing all those (a, b, c) for which system of linear equations is consistent and D be the square of the distance of the point ($-$9, 0, 0) from the plane P.
\qn{1} The value of ($|$x$|$ + 10) is \underline{\hspace{1.2cm}}
\qn{2} The value of D is \underline{\hspace{1.2cm}}
\paraq{3}{4}
Let f be a polynomial function defined only over non negative real numbers such that f(x)f(y) + 2 = f(x) + f(y) + f(xy) and f(x) is one-one function with f(0) = 1 and f $'$(1) = 2 then
\qn{3} The value of f(5) is
\qn{4} If h(x) = min$\left\{\dfrac2{\mathrm f(\mathrm x)},\mathrm x^2,|1-|\mathrm x||\right\}$, where x $\geqslant$ 0, then the number of points of non differentiability of h(x) is
\paraq{5}{6}
Let S$_\mathrm n$ be the sum of first n terms of the sequence \{a$_\mathrm n$\}. Let a$_\mathrm n$ = 5S$_\mathrm n$ + 1 holds for any n $\in$ N.\par
Let b$_\mathrm n$ = $\dfrac{4+\mathrm a_\mathrm n}{1-\mathrm a_\mathrm n}$, n $\in$ N
\qn{5} $\left(\dfrac{4^{10}\mathrm a_{10}+5}2\right)$ =
\qn{6} $\left(\dfrac{(4^{12}-1)\mathrm b_{12}-1}{2^{27}}\right)$ =
\sect{SECTION-II(ii)}
\qn{7} A ship sends five different codes A, B, C, D, E for communication, but each week uses only one of them. The code used in a definite week is randomly selected with equal chance among the four ones that have not been used in the last week. Suppose the code used in the first week is A. If the probability that A is also used in the ninth week is $\dfrac\alpha\beta$ (where $\alpha$,$\beta$ are relatively prime positive integers) then the value of $\left[\dfrac\alpha{100}\right]$ = \par
(where [.] is greatest integer function)
\qn{8} If z is a complex number such that $|$z$|$ = 1. If maximum value of $|$z$^3-$ z + 2$|$ is `a', then the sum of digit of a$^2$ is
\qn{9} If the area bounded by x = 0, y = 0, x = 2, y = 2, y $\leqslant$ e$^\mathrm x$ and y $\geqslant$ ln x is (a $-$ b ln2) sq. unit, where a $\in$ N, b $\in$ N then value of $\left(\dfrac{\mathrm a+\mathrm b}2\right)$ is
\raggedcolumns
\papersec{2}
\partsec{1}{PHYSICS}
\sect{SECTION-I(i)}
\qn{1} Two plane electromagnetic waves with the same amplitude E$_0$ and frequency $\omega$ travel in opposite directions along the z-axis, creating a standing wave. The electric field vectors of these waves are given by :\par
$\vec{\mathrm E}_1$(z, t) = E$_0$ cos(kz $-$ $\omega$t)$\hat{\mathrm i}$, $\vec{\mathrm E}_2$(z, t) = E$_0$ cos(kz + $\omega$t + $\phi$)$\hat{\mathrm i}$\par
The superposition of these two waves forms a standing wave. Select the correct statements.
\begin{qoptions}
\opt{A}{When $\phi=\dfrac\pi2$, the resulting wave is elliptically polarized and the polarization changes with position along the z-axis.}
\opt{B}{If $\phi=\pi$, the resulting wave is linearly polarized along a fixed direction in the xy-plane.}
\opt{C}{The magnetic field vector of the resulting wave oscillates perpendicular to the plane of polarization at each point.}
\opt{D}{For $\phi=\pi$, nodes of the electric field occur at positions z = $\dfrac{\mathrm n\pi}{\mathrm k}$, where n is an integer.}
\end{qoptions}
\qn{2} A small bead with mass m and charge +q can slide without friction along a non-conducting ring with radius R. In the plane of the ring, at a distance h from the center of the ring (h $<$ R), a charge +Q is fixed. Initially (see figure), the velocity of the bead is v$_0$.\par
\qfig{p2_phy_s1_q02_1}
\begin{qoptions}
\opt{A}{The maximum speed reached by the bead during the motion is $\sqrt{\mathrm v_0^2+\dfrac{2\mathrm{kqQ}}{\mathrm m}\left(\dfrac1{\sqrt{\mathrm R^2+\mathrm h^2}}-\dfrac1{\mathrm R+\mathrm h}\right)}$}
\opt{B}{The normal force acting on the bead when the speed is maximum, is zero.}
\opt{C}{The initial velocity required to complete a full revolution around the ring is\par $\sqrt{\dfrac{2\mathrm{kqQ}}{\mathrm m}\left(\dfrac1{\mathrm R-\mathrm h}-\dfrac1{\sqrt{\mathrm R^2+\mathrm h^2}}\right)}$.}
\opt{D}{The initial velocity required to complete a full revolution around the ring is\par $\sqrt{\dfrac{2\mathrm{kqQ}}{\mathrm m}\left(\dfrac1{\mathrm R-\mathrm h}+\dfrac1{\sqrt{\mathrm R^2+\mathrm h^2}}\right)}$}
\end{qoptions}
\qn{3} A thin homogeneous square plate with side length a and mass m has its center of mass at its geometric center. The plate is divided into four equal parts A,B,C,D \_ See Fig.\par
\qfig{p2_phy_s1_q03_1}
\begin{qoptions}
\opt{A}{The co-ordinates of center of mass of system after removing D is $\left(\dfrac{5\mathrm a}{12},\dfrac{\mathrm a}{12}\right)$.}
\opt{B}{The co-ordinates of center of mass of system after removing D and placing it on part A is $\left(\dfrac{3\mathrm a}8,\dfrac{\mathrm a}8\right)$.}
\opt{C}{The co-ordinates of center of mass of system after removing D and placing it on part A is $\left(\dfrac{3\mathrm a}8,\dfrac{3\mathrm a}8\right)$.}
\opt{D}{The co-ordinates of center of mass of system after removing D and placing it on part B is $\left(\dfrac{\mathrm a}2,\dfrac{3\mathrm a}8\right)$.}
\end{qoptions}
\qn{4} If the graph of binding energy $\left(\dfrac{\mathrm{B.E.}}{\mathrm A}\right)$ per nucleon versus mass number (A) looks like as shown in figure, then using the curve choose correct options :-\par
\qfig{p2_phy_s1_q04_1}
\begin{qoptions}
\opt{A}{Fusion of two nuclei with mass number 30 and 45 into one nucleons of mass number 75 will release energy.}
\opt{B}{Fission of nuclei with mass number 80 into two nuclei of equal mass number will have Q-value 256 MeV.}
\opt{C}{Fission of nuclei with mass number 150 into two nuclei of equal number will release energy.}
\opt{D}{Fusion of two nuclei with mas number 10 and 20 into a nucleus of mass number 30 will release energy.}
\end{qoptions}
\qn{5} You have three tuning forks, A, B, and C. Fork B has a frequency of 441 Hz; when A and B are sounded together, a beat frequency of 3 Hz is heard. When B and C are sounded together, the beat frequency is 4 Hz.
\begin{qoptions}
\opt{A}{The frequency of A can be 438 Hz.}
\opt{B}{The frequency of C can be 444 Hz.}
\opt{C}{The beat frequency between A and C can be 1 Hz.}
\opt{D}{The beat frequency between A and C can be 6 Hz.}
\end{qoptions}
\qn{6} A winch is used to haul sand up from the bottom of a dry dock using a crane constructed from a very light framework of girders. The crane geometry can be seen in the diagram of figure. The crane is fixed to the horizontal dock side at points A and B. The dimensions of the sides CB, CD and BD are as shown, the girder CB is vertical, and the winch is hauling a load of weight 15W.\par
\qfig{p2_phy_s1_q06_1}
\begin{qoptions}
\opt{A}{The force exerted on D by BD in 24W.}
\opt{B}{Tension in the rod CD is 15W.}
\opt{C}{If the force on the clamp at A by ground acts vertically downwards and has a magnitude 36W, the distance AB is 2L.}
\opt{D}{None of the above.}
\end{qoptions}
\sect{SECTION-I(ii)}
\paraq{7}{8}
A uniform bar of mass m is pivoted at point O and supported at the ends by two springs, as shown in figure. End P of spring PQ is subjected to a sinusoidal displacement, x(t) = x$_0$ sin $\omega$t.\par
($\ell$ = 1 m, k = 1000 N/m, m = 10 kg, x$_0$ = 1 cm, and $\omega$ = 10 rad/s.)\par
Case-I : No damper is present.\par
Case-II : A damper having damping constant C = $\dfrac{1000\sqrt2}3$ N-s/m is attached. (Neglect gravity)\par
\qfig{p2_phy_s2_para7_8_1}
\qfig{p2_phy_s2_para7_8_2}
\qn{7} The angular displacement of bar in case-I is
\begin{qoptions}
\opt{A}{0.01565 sin (10t) rad}
\opt{B}{0.01565 sin (10t + $\pi$) rad}
\opt{C}{0.01565 sin (10t + $\dfrac\pi2$) rad}
\opt{D}{0.0313 sin (10t + $\pi$) rad}
\end{qoptions}
\qn{8} The angular amplitude of bar in case-II is
\begin{qoptions}
\opt{A}{0.008 rad}
\opt{B}{0.0131 rad}
\opt{C}{0.0262 rad}
\opt{D}{0.028 rad}
\end{qoptions}
\paraq{9}{10}
Consider a one-dimensional standing electromagnetic wave in the form of E(x) = A sin(k$_\mathrm x$x) along the x-direction confined within the space between x = 0 and x = a. The wave must vanish at these two end points.
\qn{9} Find the allowed values of k$_\mathrm x$, (n = 1, 2, 3\ldots\ldots)
\begin{qoptions}
\opt{A}{k$_\mathrm x$ = $\dfrac{\mathrm n\pi}{\mathrm a}$}
\opt{B}{k$_\mathrm x$ = $\dfrac{2\mathrm n\pi}{\mathrm a}$}
\opt{C}{k$_\mathrm x$ = $\dfrac{\mathrm n\pi}{2\mathrm a}$}
\opt{D}{k$_\mathrm x$ = $\dfrac{4\mathrm n\pi}{\mathrm a}$}
\end{qoptions}
\qn{10} The String Theory predicts that our space is more than three-dimension, and the additional hidden dimensions are folded up like the dimension y on the surface of a thin cylinder shown in the figure. Suppose the radius of the cylinder is b ($<<$ a), and the electromagnetic wave on the surface now takes the form E(x,y) = A sin(k$_\mathrm x$x)cos(k$_\mathrm y$y), where y is the coordinate of the folded space around the cylinder. Find the allowed values of k$_\mathrm y$. (m = 1, 2, 3\ldots\ldots..)\par
\qfig{p2_phy_s2_q10_1}
\begin{qoptions}
\opt{A}{k$_\mathrm y$ = $\dfrac{\mathrm m}{\mathrm b}$}
\opt{B}{k$_\mathrm y$ = $\dfrac{\mathrm m}{2\mathrm b}$}
\opt{C}{k$_\mathrm y$ = $\dfrac{\mathrm m}{\pi\mathrm b}$}
\opt{D}{k$_\mathrm y$ = $\dfrac{2\mathrm m}{\pi\mathrm b}$}
\end{qoptions}
\sect{SECTION-II(i)}
\paraq{1}{2}
In the window of the physics laboratory, an inverted image of the neighboring school was photographed, created by a spherical glass flask with a radius of R = 60 mm filled with water. Take $\mu_{\mathrm{water}}$ = $\dfrac43$.\par
\qfig{p2_phy_s3_para1_2_1}
\qn{1} At what distance d$_1$ (in cm) from the center of the flask is the center of this image located?
\qn{2} Since the camera is focused on the image of the building created by the flask, the actual school building is captured only as a blur. However, you can quite accurately compare the viewing angle under which you see the window of the building from the location where the center of the lens is located with the viewing angle under which you see, from the same location, the image of the window created by the flask. In the photograph, the height of the window appears as H = 50 mm, and the height of the image of the window created by the flask as h = 6 mm. What is the distance d$_2$ (in cm) from the center of the lens to the center of the image created by the flask ?
\paraq{3}{4}
Consider two planets, A and B, of masses M$_\mathrm A$ = 1.0 $\times$ 10$^{24}$ kg and M$_\mathrm B$ = 2.0 $\times$ 10$^{24}$ kg. in circular orbits of radius R$_\mathrm A$ = 1.0 $\times$ 10$^{11}$ m and R$_\mathrm B$ = 4 $\times$ 10$^{11}$ m around a common star of mass M$_\mathrm S$ = 24 $\times$ 10$^{31}$ kg. The two planets are orbiting in the same direction. You may assume that the planets are point masses and have no rotation about their own individual axes. (Take : G = $\dfrac{20}3$ $\times$ 10$^{-11}$ N-m$^2$/kg$^2$)
\qn{3} The planets have angular momentum L$_\mathrm A$ and L$_\mathrm B$ about the star. The value of $\dfrac{\mathrm L_\mathrm A}{\mathrm L_\mathrm B}$ is :
\qn{4} The inhabitants of the planets have constructed a light indestructible link bridge, which they connect when planets A and B are closest to each other. This joins the two planets instantaneously. Determine the angular velocity $\omega$ (in rad/year) of the two planets about their centre of mass at the instant just after the connection is made.\par
\qfig{p2_phy_s3_q04_1}
\paraq{5}{6}
Consider a non-conducting sphere of radius R with centre O. A spherical cavity of radius $\dfrac R2$ is created with it's centre $\dfrac R2$ from O. The remaining sphere has charge Q uniformly distributed in it. The cavity is now filled with equal negative charge $-$Q, uniformly. [Take Q = 14 $\mu$C, R = 1cm.]\par
\qfig{p2_phy_s3_para5_6_1}
\qn{5} Find magnitude of potential energy of interation between positive and negative charges.
\qn{6} Find the net force on interaction between positive and negative charge (in kilo Newton)
\sect{SECTION-II(ii)}
\qn{7} Initially, a switch S open in the circuit shown in the figure, a capacitor of capacitance 2C carries the electric charge q$_0$, a capacitor of capacitance C is uncharged, and there are no electric currents in both coils of inductance L and 2L, respectively. The capacitor starts to discharge and at the moment when the current in the coils reaches its maximum value, the switch S is instantly shorted. If the maximum current I$_{\max}$ through the switch S thereafter is $\dfrac{\mathrm q_0}{\sqrt{\alpha\mathrm{LC}}}$, the value of $\alpha$ is\par
\qfig{p2_phy_s4_q07_1}
\qn{8} In the production of polyethylene film, a wide strip is drawn over rollers at a speed of v = 45m/s as shown in the figure. During the processing (mainly due to friction), the surface of the film acquires a uniformly distributed charge. Find the maximum values of magnetic field induction B$_{\max}$ near the surface of the film, taking into account that at an electric field strength E = 20 kV/cm, an electrical discharge occurs in the air. The value of B$_{\max}$ is $\alpha$ $\times$ 10$^{-9}$ T. The value of $\alpha$ is\par
\qfig{p2_phy_s4_q08_1}
\qn{9} A certain planet of radius R is composed of a uniform material that, through radioactive decay, generates a net power P. This results in a temperature difference between the inside and outside of the planet as heat is transferred from the interior to the surface. The thermal conductivity of the material is k. The temperature difference between the surface of the planet and the center of the planet is equal to $\dfrac{\mathrm P}{\alpha\pi\mathrm{kR}}$. The value of $\alpha$ is
\raggedcolumns
\papersec{2}
\partsec{2}{CHEMISTRY}
\sect{SECTION-I(i)}
\qn{1} Select the correct major products
\begin{qoptions}
\opt{A}{\qfigi{p2_chem_s1_q01_A}\qfigi{p2_chem_s1_q01_2}}
\opt{B}{\qfigi{p2_chem_s1_q01_B}\qfigi{p2_chem_s1_q01_4}}
\opt{C}{\qfigi{p2_chem_s1_q01_5}$\xrightarrow[\mathrm{CH_2Cl_2}]{\mathrm{mCPBA}}$\qfigi{p2_chem_s1_q01_6}}
\opt{D}{\qfigi{p2_chem_s1_q01_D}}
\end{qoptions}
\qfullbegin
\qn{2} Which of the following option(s) have one of the product(s) obtained during this reaction sequence?\par
\qfig{p2_chem_s1_q02_1}
\qfullsplit
\begin{qoptions}
\opt{A}{\qfigi{p2_chem_s1_q02_A}}
\opt{B}{\qfigi{p2_chem_s1_q02_B}}
\opt{C}{\qfigi{p2_chem_s1_q02_C}}
\opt{D}{\qfigi{p2_chem_s1_q02_D}}
\end{qoptions}
\qfullend
\qn{3} Select the correct option(s) regarding complex K$_2$Zn$_3$[Fe(CN)$_6$]$_2$ :
\begin{qoptions}
\opt{A}{It is insoluble in water but dissolves in exc. NaOH.}
\opt{B}{Complex species is paramagnetic.}
\opt{C}{Central metal uses d$^2$sp$^3$ hybrid orbitals for bond formation.}
\opt{D}{Metal ligand linkage has $\pi$-character.}
\end{qoptions}
\qn{4} Select the correct statement(s) :\par
\qfig{p2_chem_s1_q04_1}
\begin{qoptions}
\opt{A}{`Y' is see-saw shaped.}
\opt{B}{Central atom uses total three `d-orbitals' for bonding in `X'.}
\opt{C}{`Z' is planar but `X' is non-planar.}
\opt{D}{Z is white hygroscopic solid.}
\end{qoptions}
\qn{5} A metal cation M$^{+2}$ gives black precipitate with acidified H$_2$S which dissolves in hot conc. HNO$_3$. Solution of M$^{+2}$ gives brown precipitate on treatment with Potassium ferrocyanide. Select the correct statement(s) :
\begin{qoptions}
\opt{A}{M$^{+2}$(aq.) on reaction with exc. KCN produces a gas.}
\opt{B}{Metal `M' in its extraction process from its most common sulphide ore is obtained from auto-reaction process.}
\opt{C}{E$^\circ_{\mathrm M^{+2}/\mathrm M}$ is positive.}
\opt{D}{M$^{2+}_{\mathrm{(aq.)}}$ produce dark blue solution with excess Cl$^-$}
\end{qoptions}
\qn{6} pH of which of the following solution in close to 5 ? (pka of CH$_3$COOH = 5, K$_\mathrm a$ of NH$_4^+$ = 5.6 $\times$ 10$^{-10}$)
\begin{qoptions}
\opt{A}{50 mL of 0.01 M CH$_3$COOH + 50 mL of 0.01 M CH$_3$COONa}
\opt{B}{0.1 M solution of NH$_4$Cl}
\opt{C}{0.1 M weak monoprotic acid having K$_\mathrm a$ = 1 $\times$ 10$^{-9}$}
\opt{D}{5 $\times$ 10$^{-6}$ M Ba(OH)$_2$}
\end{qoptions}
\sect{SECTION-I(ii)}
\paraq{7}{8}
1g of methane diffused in 20 seconds under certain conditions. Under the same conditions, $\sqrt{20}$g of an open chain hydrocarbon (A) diffused in 40 seconds. A 10 mg of sample of (A) consumes 8.40 ml of H$_2$ gas measured at 0$^\circ$C and 760 mm pressure.\par
One of the open chain isomers of (A) on ozonolysis gives only formaldehyde \& glyoxal.
\qn{7} A six membered stable monocyclic isomer of (A) on reductive ozonolysis can NOT give which of the following products
\begin{qoptions}
\opt{A}{OHC $-$ CHO}
\opt{B}{OHC$-$CH$_2-$CHO}
\opt{C}{OHC$-$CH$_2-$CH$_2-$CHO}
\opt{D}{OHC$-$CH$_2-$CH$_2-$CH$_2-$CHO}
\end{qoptions}
\qn{8} Identify the incorrect statement.
\begin{qoptions}
\opt{A}{One of the isomers of (A) exhibits geometrical isomerism}
\opt{B}{One of the isomers of (A) gives white precipitate with Tollen 's reagent}
\opt{C}{One of the isomers of (A) is optically active.}
\opt{D}{One of the isomers of (A) on reductive ozonolysis gives only acetone.}
\end{qoptions}
\paraq{9}{10}
Consider the following series of reactions\par
\qfig{p2_chem_s2_para9_10_1}
B$_{\text{(gas)}}$ + H$_2$(g)$_{\text{excess}}\to$ C(gas)\par
C(g) + bleaching powder solution $\to$ B(g) + D + H$_2$O
\qn{9} Select the correct statement's.\par
(i) B is a diamagnetic gas having $\pi$ bond order equal to 2.\par
(ii) In molecular orbital diagram of (B) gas H.O.M.O. contains 1 nodal plane.\par
(iii) Compound C is a weak acid.\par
(iv) Compound `D' is obtained as byproduct of Solvay process.
\begin{qoptions}
\opt{A}{only (iv)}
\opt{B}{(i), (iii) and (iv)}
\opt{C}{(i) and (iv)}
\opt{D}{(i), (ii) and (iii)}
\end{qoptions}
\qn{10} Compound `D' can be used to dry :-
\begin{qoptions}
\opt{A}{NH$_3$}
\opt{B}{C$_2$H$_5$OH}
\opt{C}{Both (A) and (B)}
\opt{D}{Neither `A' nor `B'}
\end{qoptions}
\sect{SECTION-II(i)}
\paraq{1}{2}
An acid A which is an important constituent of vineagar, on reaction with red P and Br$_2$ gives a monobromoderivative (B) which on reaction with NH$_3$ gives white solid (C). However ethanal on reaction with a mixture of ammonium chloride and sodium cyanide undergoes Strecker synthesis (synthesis of amino acids by the reaction of an aldehyde with cyanide in presence of ammonia) to give a product which on acidic hydrolysis gives another high melting solid (D). Let x and y are the total number of optical isomers of (C) and (D) respectively while z is the total number of dipeptides (excluding stereoisomers) that (C) and (D) can form.
\qn{1} The value of x + y is \underline{\hspace{1.5cm}}.
\qn{2} The value of z is \underline{\hspace{1.5cm}}.
\paraq{3}{4}
Consider the following sequence of reactions.\par
(i) PbO$_2$ + conc. HNO$_3\to$ Pb(NO$_3$)$_2$ + H$_2$O + \underline{X} (g)\par
(ii) 2NaHCO$_3\xrightarrow{\Delta}$ Na$_2$CO$_3$ + H$_2$O + \underline{Y} (g)\par
(iii) FeS$_{2(8)}$ + X(g) $\to$ Fe$_2$O$_3$ + \underline{Z} (g)
\qn{3} Assume ideal behavior of all three gases, 0.5 moles of each X and Y gas are taken in a container and expanded reversibly and adiabatically from 1 dm$^3$ to 4 dm$^3$ starting from 27$^\circ$C. Find final temperature (K).\par
Given (4)$^{0.4}$ = 1.74
\qn{4} If equal moles of X,Y,Z (assuming ideal behavior of the gas) are expanded adiabatically and reversibly from the same initial state to same final volume. The magnitude of work is maximum for a gas. What is the value of $\gamma$ for that gas, if all the degree of freedom are active.
\paraq{5}{6}
The following reacations are occurring in solvay process used for the manufacture of sodium carbonate.\par
(A) $\xrightarrow{\Delta}$ (B) + CO$_2$\par
(B) + H$_2$O $\to$ (C)\par
\qfig{p2_chem_s3_para5_6_1}
(E) + CO$_2$(excess) $\to$ (F)\par
(F) + NaCl $\to$ (G) $\xrightarrow{\Delta}$ sodium carbonate.\par
\textbf{(Given that Atomic Mass : Na = 23, K = 39, Ca = 40, Mg = 24, Fe = 56, Cu = 63.5, O = 16, S = 32, H = 1, C = 12, Cl = 35.5, F = 19, Zn = 65)}
\qn{5} If the molar mass of G is X g mole$^{-1}$ and that of (D) is Y g mole$^{-1}$. Then (X + Y) is\underline{\hspace{1cm}}
\qn{6} The number of molecules of F are obtained on reaction of $\left(\dfrac12\right)$ mole of (E) with $\left(\dfrac14\right)$ mole of CO$_2$ is Q $\times$ 10$^{21}$. What is the value of Q ?
\sect{SECTION-II(ii)}
\qn{7} In how many steps decarboxylation reaction is taking place.\par
\qfigi{p2_chem_s4_q07_1}$\xrightarrow[\text{Step--1}]{\mathrm{H_3O^+}}\xrightarrow[\text{Step--2}]{\Delta}\xrightarrow[\text{Step--3}]{\mathrm{I_2/OH^-}}$\par$\xrightarrow[\text{Step--4}]{\mathrm{KMnO_4/H^+}}\xrightarrow[\text{Step--5}]{\Delta}\xrightarrow[\text{Step--6}]{\mathrm{NaOH/CaO/\Delta}}$
\qn{8} Two mole of non-volatile weak electrolyte XY was dissolved in 20 dm$^3$ of water. If boiling point of the solution was found to be 100.08$^\circ$C and K$_\mathrm b$ of H$_2$O is 0.5, then calculate the conductivity of the solution in S m$^{-1}$ unit. Assuming molality is equal to molarity.\par
$\Lambda^\circ_\mathrm m$ of XY = 500 Scm$^2$/mole
\qn{9} 200 cm$^3$, 0.2 M BaCl$_2$(aq.) is mixed with 0.1M, 600 cm$^3$ Na$_2$SO$_4$ solution. Find osmotic pressure of resulting solution (in atm). Round off up to nearest integer value.\par
R = 0.0821 L-atm K$^{-1}$ mol$^{-1}$\par
T = 37$^\circ$C
\raggedcolumns
\papersec{2}
\partsec{3}{MATHEMATICS}
\sect{SECTION-I(i)}
\qn{1} Let f(x) be a polynomial function satisfying 0 $<$ xf (y) $<$ yf (x) $\forall$ x, y such that 0 $<$ x $<$ y $<$ 1 and f (0) = 0 then :
\begin{qoptions}
\opt{A}{f$'$(x) $<$ f (1)}
\opt{B}{f(1) $<$ 2 $\displaystyle\int_0^1$f(x) dx}
\opt{C}{3f $\left(\dfrac13\right)$ $>$ 2 f $\left(\dfrac12\right)$}
\opt{D}{6f $\left(\dfrac16\right)$ $<$ 5 f $\left(\dfrac15\right)$}
\end{qoptions}
\qn{2} Let a, b , c be distinct complex numbers with $|$a$|$ = $|$b$|$ = $|$c$|$ = 1 and z$_1$, z$_2$ be the roots of the equation az$^2$ + bz + c = 0 with $|$z$_1|$ = 1. Also P and Q are the points representing the complex number z$_1$ and z$_2$ respectively in the complex plane with $\angle$POQ = $\theta$ ( where O being the origin) then which of the following is/are correct?
\begin{qoptions}
\opt{A}{b$^2$ = ac}
\opt{B}{$\theta=\dfrac{2\pi}3$}
\opt{C}{PQ = $\sqrt3$}
\opt{D}{$|$z$_1$ + z$_2|$ = 1}
\end{qoptions}
\qn{3} Let $<$T$_\mathrm n>$ be a sequence such that T$_\mathrm n^3$ + 2T$_\mathrm n$ = T$_{\mathrm n+1}$ $\forall$ n $\in$ N, and T$_1$ = 1, then :
\begin{qoptions}
\opt{A}{$\displaystyle\sum_{\mathrm n=1}^{100}(\mathrm T_\mathrm n^3+\mathrm T_\mathrm n+1)=\mathrm T_{101}+99$}
\opt{B}{$\displaystyle\sum_{\mathrm n=1}^{100}(\mathrm T_\mathrm n^3+\mathrm T_\mathrm n+1)=\mathrm T_{101}+100$}
\opt{C}{$\displaystyle\prod_{\mathrm n=1}^{100}\left(\mathrm T_\mathrm n^2+2\right)=\mathrm T_{100}$}
\opt{D}{$\displaystyle\prod_{\mathrm n=1}^{100}\left(\dfrac{\mathrm T_{\mathrm n+1}}{\mathrm T_\mathrm n}-\mathrm T_\mathrm n^2\right)=2^{100}$}
\end{qoptions}
\qn{4} Let straight line y = mx + 4 meets the curve 3x$^2-$ (1 $-$ 3a) xy $-$ ay$^2$ = 0 at two points A and B such that $\angle$AOB = 90$^\circ$ $\forall$ m $\in$ R $-$ \{m$_1$,m$_2$\} where m$_1$ $<$ m$_2$ and `O' is the origin. Identify which of the following statements is/are correct ?
\begin{qoptions}
\opt{A}{m$_1$ + m$_2$ = $\dfrac{10}3$}
\opt{B}{am$_1$ + m$_2$ = 2}
\opt{C}{If m = 2, then area of $\Delta$AOB = $\dfrac{80}7$ sq. units}
\opt{D}{If m = 2, then area of$\Delta$AOB = $\dfrac{85}7$ sq. units}
\end{qoptions}
\qn{5} Let A = $\begin{bmatrix}\mathrm a&\mathrm b&\mathrm c\\\mathrm b&\mathrm c&\mathrm a\\\mathrm c&\mathrm a&\mathrm b\end{bmatrix}$ and B = $\begin{bmatrix}\mathrm{bc}-\mathrm a^2&\mathrm{ca}-\mathrm b^2&\mathrm{ab}-\mathrm c^2\\\mathrm{ca}-\mathrm b^2&\mathrm{ab}-\mathrm c^2&\mathrm{bc}-\mathrm a^2\\\mathrm{ab}-\mathrm c^2&\mathrm{bc}-\mathrm a^2&\mathrm{ac}-\mathrm b^2\end{bmatrix}$ be two non-singular matrices such that (A$^2-$ 2I) B = O where a $>$ b $>$ c $>$ 0, then which of the following statements is/are correct?\par
[\textbf{Note:} I is an identify matrix of order 3 and Tr. (P) and det. (P) denote trace and value of the determinant of square matrix P respectively.]
\begin{qoptions}
\opt{A}{Tr. (AB) = 6$\sqrt2$}
\opt{B}{Tr. (AB) = $-$6$\sqrt2$}
\opt{C}{det. (A $-$ $\sqrt2$B) = 54$\sqrt2$}
\opt{D}{det. (A $-$ $\sqrt2$B) = $-$54$\sqrt2$}
\end{qoptions}
\qn{6} If $\lambda_1>\lambda_2>\lambda_3>\ldots>\lambda_\mathrm n$ be the values of `$\lambda$' for which equation x$^3-$ 3 x + $\lambda$ = 0 has all integral roots, if h(x) = x$^3-$ 3x + $|\lambda_1|$ + $|\lambda_2|$\ldots..+ $|\lambda_\mathrm n|$ for x $\in$ [1, $\infty$) and g(x) be the inverse of h(x) then:
\begin{qoptions}
\opt{A}{($\lambda_1$)$^2$ + ($\lambda_2$)$^2$+\ldots..+($\lambda_\mathrm n$)$^2$ = a, number which is not the perfect square of an integer.}
\opt{B}{Area bounded by g(x) and y-axis between limits y = 2 to y = 5 is $\dfrac{531}4$}
\opt{C}{$\displaystyle\int_1^3$(x$^3-$3x)dx + $\displaystyle\int_2^{22}$g(x) dx = 64}
\opt{D}{$\displaystyle\int_1^3$(x$^3-$3x)dx + $\displaystyle\int_2^{22}$g(x) dx = 56}
\end{qoptions}
\sect{SECTION-I(ii)}
\paraq{7}{8}
Consider, f(x) = $\displaystyle\lim_{\mathrm n\to\infty}\dfrac{\mathrm{sgn}(\sqrt{\mathrm{ac}}-\mathrm b)\mathrm e^{\mathrm{nx}}+\mathrm x^2+\mathrm f}{2\mathrm e^{\mathrm{nx}+\mathrm x}+\mathrm x+\mathrm d}$ where a $>$ b $>$ c $>$ 0 and d, f $\in$ R\par
[Note: sgn (y) denotes the signum function of y.]
\qn{7} If a, b and c are in A.P. and f (x) is continuous for all x $\in$ R, then the value of (2f + d + 1) is equal to :
\begin{qoptions}
\opt{A}{0}
\opt{B}{1}
\opt{C}{$-$1}
\opt{D}{$\dfrac{-1}2$}
\end{qoptions}
\qn{8} If a, b and c are in G.P. and f(x) is continuous for all x $\in$ R (d $<$ 0), then number of solutions of the equation f(x) = $||$x $-$ 4$|-$ 2$|-$ 1 is (are) :
\begin{qoptions}
\opt{A}{0}
\opt{B}{2}
\opt{C}{3}
\opt{D}{4}
\end{qoptions}
\paraq{9}{10}
Let B$_\mathrm n$ denotes the event that n fair dice are rolled once with P(B$_\mathrm n$) = $\dfrac1{2^\mathrm n}$, n $\in$ N, e.g.\par
P(B$_1$) = $\dfrac12$, P(B$_2$) = $\dfrac1{2^2}$, P(B$_3$) = $\dfrac1{2^3}$, \ldots\ and P(B$_\mathrm n$) = $\dfrac1{2^\mathrm n}$. Hence B$_1$, B$_2$, B$_3$,\ldots,B$_\mathrm n$ are pair wise mutually exclusive and exhaustive events as n $\to\infty$. The events A occurs with atleast one of the event B$_1$, B$_2$, B$_3$,\ldots\ldots\ldots.B$_\mathrm n$ and denotes that the sum of the numbers appearing on the dice is S.
\qn{9} If even number of dice has been rolled, the probability that S = 4, is
\begin{qoptions}
\opt{A}{very closed to $\dfrac12$}
\opt{B}{very closed to $\dfrac14$}
\opt{C}{very closed to $\dfrac18$}
\opt{D}{very closed to $\dfrac1{16}$}
\end{qoptions}
\qn{10} Probability that greatest number on the dice is 4 if three dice are known to have been rolled, is
\begin{qoptions}
\opt{A}{$\dfrac{37}{216}$}
\opt{B}{$\dfrac{64}{216}$}
\opt{C}{$\dfrac{27}{216}$}
\opt{D}{$\dfrac{31}{216}$}
\end{qoptions}
\sect{SECTION-II(i)}
\paraq{1}{2}
An ellipse whose major axis is parallel to x-axis such that segment of a focal chord are 1 and 2 units. The line $\dfrac xa+\dfrac yb+\dfrac1c=0$ are the chords of the ellipse such that a, c, b are in H.P. and are bisected by the point at which they are concurrent, the equation of auxiliary circle is 2x$^2$ + 2y$^2$ + 4px + 4qy $-$ 16p + 1 = 0 then
\qn{1} Find p + q
\qn{2} If length of latus rectum of ellipse is L, then find [L]. (where [.] represent Greatest integer function)
\paraq{3}{4}
If n be a natural number define polynomial f$_\mathrm n$ (x) of n$^{\mathrm{th}}$ degree as follows f$_\mathrm n$ (cos$\theta$) = cos n$\theta$ i.e. f$_2$ (x) = 2x$^2-$1 ; f$_3$ (x) = 4x$^3-$ 3x. Then.
\qn{3} $\left(\mathrm x+\sqrt{\mathrm x^2-1}\right)^{10}+\left(\mathrm x-\sqrt{\mathrm x^2-1}\right)^{10}$ is equal to $\lambda$ f$_\mathrm k$ (x). Then find value of $\lambda$ + k.
\qn{4} f$_6$ (x) is equal to px$^6$ + qx$^4$ + rx$^2$ + s then find value of 7 (p + q + r $-$ s)
\paraq{5}{6}
Let P$_1$ : x + y + 2z $-$ 3 = 0 and P$_2$ : x $-$ 2y + z = 4 be two planes. Let A(1,3,4) and B(3,2,7) be two points in R$^3$.
\qn{5} The equation of plane P$_3$ through line of intersection of P$_1$ and P$_2$ and length of projection upon which of the line segment AB is greatest is equal to ax + by + cz + 1 = 0, then find the value of a + b + c. (where a,b,c $\in$ I)
\qn{6} The equation of plane P$_4$ through the line of intersection of P$_1$ and P$_2$ and the length of projection upon which of the line segment AB is least is equal to ax + by + cz $-$ 7 = 0, then find the value of a + 2b + c. (where a,b,c $\in$ I)
\sect{SECTION-II(ii)}
\qn{7} The sum of infinite series\par
$\dfrac1{15}+\dfrac1{30}+\dfrac1{50}+\dfrac1{75}+\ldots\ldots..\infty$ is $\dfrac{\mathrm k}{20}$, then the value of k is
\qn{8} Given that for a, b, c, d $\in$ R, if a sec (200$^\circ$) $-$ c tan (200$^\circ$) = d and b sec (200$^\circ$) + d tan (200$^\circ$) = c then\par
$\dfrac{\mathrm a^2+\mathrm b^2+\mathrm c^2+\mathrm d^2}{\mathrm{ac}-\mathrm{bd}}$ = pcosec(200$^\circ$) find the value of p\underline{\hspace{1.2cm}}.
\qn{9} If I$_\mathrm n$ = $\displaystyle\int_0^\infty$e$^{-\mathrm x}$(sin x)$^\mathrm n$dx(n $>$ 1), then the value of $\dfrac{101\,\mathrm I_{10}}{\mathrm I_8}$ is equal to \underline{\hspace{2cm}}.
\end{multicols}
\end{document}
